\section{The certificate}\label{sec:certificate}


This section states the one computed ingredient of the main theorem, says in outline how the certificate
was found, and describes what it looks like. Everything in this
section except the statement of \cref{fact:H} is descriptive: the proof of \cref{thm:main}
uses the two certificate files and nothing else from this section. How the certificate is verified is the
subject of \cref{sec:verification}; the search that found it and two observations about the search
space are described in detail in \cref{sec:supp-search}.

\subsection{The computed fact}\label{ssec:fact}

Recall the abstract game of \cref{def:states} and the notion of a valid certificate $(\Tbl,\Lv,\Rmax)$ of
\cref{def:certificate}: $\Tbl$ assigns a cell (the \emph{table move})
to canonical Q-states, $\Lv$ is a set of canonical P-states (the \emph{leaves}), and validity means that every
state reachable from the root through the rules of the game is either a listed leaf satisfying
$\LEAF$, or an expanded P-state of round $t<\Rmax$ at which $\LOST$ is false, or a Q-state with a legal table
move.

\begin{fact}[The certificate is valid]\label{fact:H}
Let $\Tbl$ be the table of $93{,}164$ canonical Q-states with one move each contained in the file
\fn{results/final/strategy.tbl}, and let $\Lv$ be the set of $106{,}723$ canonical P-states listed in the file
\fn{results/final/leaves.txt}. Then $(\Tbl,\Lv,10)$ is a valid certificate in the sense of
\cref{def:certificate}. Its reachable set consists of $153{,}135$ P-states, of which $46{,}412$ are expanded
and $106{,}723$ are leaves, and of $93{,}164$ Q-states; every entry of both files is reached. In
particular, the root $\Pst(\emptyset,\emptyset,0)$ of the abstract game is \WON{} for $\Rmax=10$
\textup{(\cref{lem:certificate}\,(b))}.

\established the program \fn{verify/treecheck.cpp} (run with the option \fn{--exact}; the literal
validation of all flows, option \fn{--paranoid}, is its default; output
\fn{results/final/treecheck.json}) and, independently, the program
\fn{verify/spotcheck.py} in its exhaustive mode (output \fn{results/final/spotcheck_full.json}). Two further
re-expansions, written separately during the internal review of the project, give the same result and the
same numbers of states and of children in every round, and the self-check of the program that produced the
files gives the same totals. During an internal refereeing of the first version of this paper two more
verifiers were written from the text of the paper; both accept the certificate and reproduce the same
numbers in every round, and a
transcription of \cref{alg:verify} alone, without the leaf test, reproduces them as well. See
\cref{sec:verification} and \cref{tab:verifiers}.
\end{fact}

The two files are identified by the SHA-256 checksums printed in \cref{ssec:intro-artifact}.
They have $5{,}376{,}504$ and $6{,}445{,}022$ bytes; their formats are specified in \cref{sec:formats}.

Apart from the two standard theorems quoted in \cref{ssec:intro-nature}, \cref{fact:H} is the only
statement on which \cref{thm:main} depends that is not proved in
this paper. It is a finite statement about two explicit files, and \cref{sec:verification} describes
precisely what a program has to do in order to check it. Since $\Rmax=10$, every expanded P-state has
$t\le 9$, and by \cref{prop:lost-stones} the predicate $\LOST$ is false at such a state; so for this
certificate the condition ``$\LOST$ is false'' is automatically satisfied. The verifiers evaluate it
nevertheless.

\subsection{How the certificate was found}\label{ssec:search}

The certificate was produced by an AND--OR search (\fn{src/search.cpp}). \emph{The search is a heuristic
device and is not part of the proof.} None of its internals is trusted: its move ordering, its cut-offs, its
hash tables and its parallelism could all be wrong without affecting the theorem, because only its output is
used, and the output is verified from scratch by programs that share no code with it. The search is
described in \cref{ssec:supp-search-details} for readers who want to reproduce or extend the computation.

In outline, the search is a depth-first evaluation of the definition of \WON{} (\cref{def:won}) in which
every state is replaced by its canonical form and results are stored in a transposition table, so that
each symmetry class is treated once. At a P-state the leaf test comes first. At a Q-state at most four of
the first ten candidate moves of a heuristic order are examined in full: busy cells, ordered with the
help of the leaf table of \cref{ssec:leaftable}, and cells suggested by the witness that a failed leaf
test returns. A \emph{relaxed game}, without $\LOST$ and without the forced children, is evaluated first
and is used only to abandon hopeless states quickly. Because of these limits the search is incomplete: a
failure of the search does not show that a state is not \WON.

The run that produced the certificate took $57.9$ seconds of wall-clock time on a four-core machine
(\cref{sec:repro}) and expanded $49{,}397$ P-states, of which $46{,}412$ belong to the final tree
(\fn{results/final/stats.json}). The program is designed so that the move recorded for a state does not
depend on the order in which the threads finish; we have not proved that the output is a deterministic
function of the sources and of the options. What we observed is that a later run from scratch with two
threads instead of four reproduced both files byte for byte (\cref{sec:repro}). The recursion, the
candidate lists, the relaxed pre-pass, the command line and the statistics of the run are described in
\cref{ssec:supp-search-details}.

\subsection{Size and shape of the tree}\label{ssec:shape}

\Cref{tab:tree} lists the reachable states of the certificate by round. States are counted up to the
$1{,}152$ symmetries, i.e.\ as canonical forms. A P-state of round $t$ is a position after $t$ complete
rounds, with Maker to move; a Q-state of round $t$ is a position after the $t$-th move of Maker, in
which Breaker is to choose a table move.

\begin{table}[ht]
\centering\small
\begin{tabular}{@{}rrrrrrrr@{}}
\toprule
 & \multicolumn{3}{c}{P-states} & & \multicolumn{3}{c}{children entering round $t$}\\
\cmidrule(lr){2-4}\cmidrule(l){6-8}
$t$ & all & leaves & expanded & Q-states & threat (c) & danger (d) & \VOID{} (ii)\\
\midrule
 0 & 1 & 0 & 1 & 0 & 0 & 0 & 0\\
 1 & 2 & 0 & 2 & 2 & 0 & 0 & 0\\
 2 & 5 & 0 & 5 & 5 & 0 & 0 & 0\\
 3 & 14 & 0 & 14 & 14 & 0 & 0 & 0\\
 4 & 35 & 0 & 35 & 35 & 0 & 0 & 0\\
 5 & 181 & 0 & 181 & 101 & 824 & 0 & 0\\
 6 & 1{,}718 & 0 & 1{,}718 & 446 & 4{,}075 & 0 & 1\\
 7 & 7{,}227 & 0 & 7{,}227 & 3{,}487 & 14{,}460 & 0 & 6\\
 8 & 34{,}801 & 3{,}382 & 31{,}419 & 14{,}362 & 47{,}020 & 20{,}416 & 54\\
 9 & 93{,}124 & 87{,}314 & 5{,}810 & 62{,}381 & 72{,}352 & 6{,}060 & 88\\
10 & 16{,}027 & 16{,}027 & 0 & 12{,}331 & 4{,}767 & 1{,}799 & 2\\
\midrule
total & 153{,}135 & 106{,}723 & 46{,}412 & 93{,}164 & 143{,}498 & 28{,}275 & 151\\
\bottomrule
\end{tabular}
\caption{The reachable set of the certificate by round $t$ (\fn{results/final/treecheck.json}; identical
numbers from the programs of \cref{tab:verifiers} that report them per round). States are counted up to
symmetry. The columns
``threat'' and ``danger'' count the non-\VOID{} children of type (c) and (d) of the
expanded P-states of round $t-1$. A child is counted once for every expanded state of which it is a
child; by \cref{lem:children-distinct} this is the number of data $(l,e)$, respectively
$(\{a,b\},Q_a,Q_b)$, for which a non-\VOID{} child is generated in lines
\ref{ln:childa}--\ref{ln:childd} of \cref{alg:verify}. Children of different states that have the same
canonical form are counted separately. The last column counts the children of type (ii) of
Q-states of round $t$ that are \VOID{} and are therefore dropped.}
\label{tab:tree}
\end{table}

In all, the reachable set has $246{,}299$ states. The following features can be read off.

\begin{itemize}[leftmargin=1.6em]
\item \emph{Depth.} All leaves occur after rounds $8$, $9$ and $10$: $3{,}382$, $87{,}314$ and $16{,}027$
  of them. No P-state of round $10$ is expanded, as $\Rmax=10$ requires.
\item \emph{Forced events.} A threat needs five stones of Maker on a live line, so the first threat children
  enter in round $5$; a danger pair needs eight stones, so the first danger children enter in round $8$.
  Up to round $4$ the tree is tiny ($57$ P-states), and its later growth is caused almost entirely by
  the forced children: $143{,}498$ threat children and $28{,}275$ danger children are generated, against
  $46{,}412$ children of type (a).
\item \emph{Children that are \VOID.} Exactly $151$ generated children are \VOID. All of them are children of type
  (ii): the table move $m$ of a Q-state $\Qst(B,F,t)$ lies on a live line that already contains four
  cells of $F$, so that $F\cup\{m\}$ is not $B$-admissible. (In real play this means that $m$ cannot be
  occupied by an unrevealed stone, because Maker would then have had five stones on a live line and the
  round would have been a threat round.) No child of type (c) or (d) is \VOID. Of the $93{,}164$ Q-states,
  $54{,}482$ have $|F|<t$ and hence a child of type (ii).
\item \emph{Unrevealed stones at the leaves.} \Cref{tab:leaf-theta} splits the leaves by
  $\theta=t-|F|$, the number of stones of Maker whose position the abstract state does not know and which
  the leaf condition must therefore withstand in every admissible position. For $60.8\%$ of the leaves
  all stones of Maker are revealed ($\theta=0$), and then $\LEAF(B,F,t)$ only asks for a pairing for $(B,F)$.
  At the other extreme there is exactly one leaf with $\theta=8$: the state $\Pst(B,\emptyset,8)$ with $B$
  in the orbit of the set $\Teight$ of \cref{ex:T8}.
\item \emph{Non-busy cells.} Of the $93{,}164$ table moves, $434$ are cells that are not busy: $16$, $399$,
  $11$ and $8$ of them in rounds $7$, $8$, $9$ and $10$. All table moves of rounds $1$ to $6$ are busy
  cells. (Breaker's stones on non-busy cells also arise from forced replies, which can be any cell: only
  $15{,}794$ of the leaves have all stones of Breaker on busy cells; $84{,}556$ have one stone on a non-busy
  cell, $6{,}343$ have two, and $30$ have three.) A computation with the search program indicates that
  non-busy table moves cannot be avoided for $\Rmax\le 10$ (\cref{rem:observations}).
\end{itemize}

\begin{table}[ht]
\centering\small
\begin{tabular}{@{}rrrrrrrrrrr@{}}
\toprule
 & \multicolumn{9}{c}{$\theta=t-|F|$} & \\
\cmidrule(lr){2-10}
$t$ & 0 & 1 & 2 & 3 & 4 & 5 & 6 & 7 & 8 & total\\
\midrule
 8 & 1{,}862 & 940 & 340 & 119 & 69 & 36 & 12 & 3 & 1 & 3{,}382\\
 9 & 53{,}397 & 26{,}291 & 6{,}102 & 1{,}189 & 260 & 63 & 11 & 1 & 0 & 87{,}314\\
10 & 9{,}582 & 5{,}472 & 825 & 126 & 14 & 6 & 2 & 0 & 0 & 16{,}027\\
\midrule
total & 64{,}841 & 32{,}703 & 7{,}267 & 1{,}434 & 343 & 105 & 25 & 4 & 1 & 106{,}723\\
\bottomrule
\end{tabular}
\caption{The leaves of the certificate by round $t$ and by the number $\theta$ of unrevealed stones of
Maker (recounted from \fn{results/final/leaves.txt}).}
\label{tab:leaf-theta}
\end{table}

\subsection{The leaf table of busy sets}\label{ssec:leaftable}

Before the search, the leaf condition $\LEAF(B,\emptyset,k)$ was evaluated for every set $B$ of $k$ busy
cells, $k=8,9,10$, up to symmetry (\fn{src/leaftable.cpp}; results in
\fn{results/leaftable/phi_k8.jsonl}, \fn{phi_k9.jsonl}, \fn{phi_k10.jsonl}). The value answers the
question: if Breaker owns exactly the stones $B$ after round $k$ and nothing is known about Maker's $k$
stones except that the position is admissible, can Breaker switch to a pairing strategy? Since $F=\emptyset$
and $k\le 10$, \cref{prop:leaf-equals-C} shows that $\LEAF$ is here equivalent to the pairing condition
$\CC(B,\emptyset,k)$ itself, so the table is exact for that question. The result is in \cref{tab:leaftable}.

\begin{table}[ht]
\centering\small
\begin{tabular}{@{}rrrr@{}}
\toprule
$k$ & orbits of busy $k$-sets & of which $\slack(B)\ge k$ & of which $\LEAF(B,\emptyset,k)$\\
\midrule
 8 & 1{,}057 & 6 & 6\\
 9 & 1{,}669 & 947 & 172\\
10 & 2{,}468 & 2{,}185 & 0\\
\bottomrule
\end{tabular}
\caption{The leaf table of busy sets (computed; exact integer max-flow). By \cref{cor:slack-necessary}
the condition $\slack(B)\ge k$ is necessary for $\LEAF(B,\emptyset,k)$.}
\label{tab:leaftable}
\end{table}

For $k=8$ the six orbits with $\LEAF$ are exactly the six orbits of slack $8$, which is the largest slack
of a busy $8$-set. Their canonical representatives are
\begin{align*}
O_1&=\{(0,0,0),(0,0,5),(1,1,1),(1,4,1),(2,3,3),(3,2,2),(4,1,4),(5,5,0)\}, & |\G O_1|&=288,\\
O_2&=\{(0,0,0),(0,0,5),(1,1,1),(1,4,1),(2,3,3),(3,2,2),(4,1,4),(5,5,5)\}, & |\G O_2|&=576,\\
O_3&=\{(0,0,0),(0,0,5),(1,1,1),(1,4,1),(2,3,3),(3,2,2),(4,4,4),(5,5,5)\}, & |\G O_3|&=288,\\
O_4&=\{(0,0,0),(0,0,5),(1,1,1),(2,3,2),(2,3,3),(3,2,2),(4,4,4),(5,5,0)\}, & |\G O_4|&=576,\\
O_5&=\{(0,0,0),(0,0,5),(1,1,1),(2,3,2),(2,3,3),(3,2,2),(4,4,4),(5,5,5)\}, & |\G O_5|&=576,\\
O_6&=\{(0,0,0),(0,0,5),(1,4,1),(2,3,3),(3,2,2),(4,1,4),(5,5,0),(5,5,5)\}, & |\G O_6|&=36,
\end{align*}
where $|\G O_i|$ is the size of the orbit; together the six orbits contain $2{,}340$ sets. Each $O_i$
leaves exactly $100$ of the $148$ lines alive. The set $O_6$ is the canonical form of
\[
\Teight=\{(0,0,0),(5,5,5),\,(1,1,4),(4,4,1),\,(2,3,2),(3,2,3),\,(5,0,0),(0,5,5)\}
\]
of \cref{ex:T8}; it is the only one of the six that consists of four antipodal pairs (one pair on each
space diagonal), the other five contain three antipodal pairs. For $k=9$ the $172$ orbits with $\LEAF$
contain $83{,}904$ sets. For $k=10$ no busy set qualifies. This is proved in \cref{prop:busy10}: every
set of ten busy cells leaves a triangle of live lines, which can absorb $4+3+3=10$ admissible stones of
Maker; the computation confirms it ($0$ of the $2{,}468$ orbits satisfy $\LEAF$).

The leaf table is not used by the proof. It guides the move ordering of the search (the ``good'' sets
of \cref{ssec:supp-search-details}), it explains the shape of the tree (the strategy steers toward one of the six
orbits when this is possible, and the earliest leaves appear after round $8$), and its values serve as
test vectors (\cref{ssec:fmt-orbits}). The numbers of orbits ($527$ for $k=7$ and the three values in the
table) were obtained both by Burnside's lemma and by direct enumeration.

\subsection{Two observations about the search space}\label{ssec:observations}

\begin{remark}[Two observations made with the search program alone]\label{rem:observations}
Two features of the certificate are explained by runs of the search program. These runs have not been
verified by an independent program and are not used anywhere in the proof; the details are given in
\cref{ssec:supp-observations}.

(i) \emph{Busy cells alone seem not to suffice for $\Rmax\le 10$.} When the table moves are restricted to
the $24$ busy cells, the search program reports that the root is not \WON{} for $\Rmax=10$, and hence for
no smaller $\Rmax$. The obstruction is the set of the $14$ lines inside a face of the cube, which only
stones inside that face can kill. The conclusion rests on the correctness of the search program in this
mode and on a pruning bound that is not proved in this paper. This is why the search was run with
non-busy moves allowed and why $434$ table moves of the certificate are non-busy cells.

(ii) \emph{$\Rmax=9$ was not achieved.} With $\Rmax=9$ the search fails, but with truncated lists of
candidate moves, so that the failure does not show that the root is not \WON{} for $\Rmax=9$. Whether
$\Rmax=9$ suffices is open; the theorem does not depend on it.
\end{remark}

\subsection{The first rounds of the strategy}\label{ssec:walk}

To give an impression of what the certified strategy does, \cref{tab:opening} lists all table entries of
rounds $1$ to $3$, translated from cell ids to coordinates. Each entry is written in the canonical frame of
its own state, so the frame changes from entry to entry; in real play the move is translated back with the
symmetry that carries the real position to its canonical form (\cref{def:sigma}).

\begin{table}[ht]
\centering\footnotesize
\begin{tabular}[t]{@{}clll@{}}
\toprule
$t$ & $B$ & $F$ & move\\
\midrule
1 & $\emptyset$ & $\emptyset$ & $(0,0,0)$\\
1 & $\emptyset$ & $(0,0,0)$ & $(1,1,1)$\\
\midrule
2 & $(0,0,0)$ & $\emptyset$ & $(5,5,5)$\\
2 & $(0,0,0)$ & $(1,1,1)$ & $(5,5,5)$\\
2 & $(0,0,0)$ & $(5,5,5)$ & $(1,1,4)$\\
2 & $(0,0,0)$ & $(1,1,1),(5,5,5)$ & $(1,1,4)$\\
2 & $(0,0,0)$ & $(1,1,4),(5,5,5)$ & $(1,4,4)$\\
\midrule
3 & $(0,0,0),(5,5,5)$ & $\emptyset$ & $(1,1,4)$\\
3 & $(0,0,0),(5,5,5)$ & $(1,1,1)$ & $(1,1,4)$\\
3 & $(0,0,0),(5,5,5)$ & $(1,1,4)$ & $(2,2,3)$\\
3 & $(0,0,0),(5,5,5)$ & $(1,1,1),(1,1,4)$ & $(1,4,1)$\\
\bottomrule
\end{tabular}\hfill
\begin{tabular}[t]{@{}clll@{}}
\toprule
$t$ & $B$ & $F$ & move\\
\midrule
3 & $(0,0,0),(5,5,5)$ & $(1,1,4),(2,2,3)$ & $(1,4,4)$\\
3 & $(0,0,0),(5,5,5)$ & $(1,1,1),(1,1,4),(1,4,1)$ & $(1,4,4)$\\
3 & $(0,0,0),(5,5,5)$ & $(1,1,4),(1,4,4),(2,2,3)$ & $(1,4,1)$\\
3 & $(0,0,0),(1,1,4)$ & $(4,4,1)$ & $(5,5,5)$\\
3 & $(0,0,0),(1,1,4)$ & $(0,0,5),(4,4,1)$ & $(5,5,5)$\\
3 & $(0,0,0),(1,1,4)$ & $(0,5,0),(4,4,1)$ & $(5,5,5)$\\
3 & $(0,0,0),(1,1,4)$ & $(4,4,1),(5,5,5)$ & $(0,5,5)$\\
3 & $(0,0,0),(1,1,4)$ & $(0,0,5),(4,4,1),(5,5,5)$ & $(0,5,5)$\\
3 & $(0,0,0),(1,1,4)$ & $(0,5,0),(4,4,1),(5,5,5)$ & $(0,5,5)$\\
3 & $(0,0,0),(1,1,4)$ & $(0,5,5),(4,4,1),(5,5,5)$ & $(5,0,5)$\\
\bottomrule
\end{tabular}
\caption{All $2+5+14$ entries of the strategy table for rounds $t=1,2,3$, in coordinates: Breaker's stones
$B$, the revealed stones $F$ of Maker, and the table move. All these table moves are busy cells.}
\label{tab:opening}
\end{table}

The pattern in these entries is simple. Breaker opens in a corner. As long as nothing is revealed,
Breaker completes antipodal pairs of busy cells: the answer to $B=\{(0,0,0)\}$ is the antipode $(5,5,5)$, and in
round $4$ the entry for $B=\{(0,0,0),(1,1,4),(4,4,1)\}$, $F=\emptyset$ is again $(5,5,5)$. When the wanted
cell is occupied, the stone found there is revealed and Breaker turns to another busy cell: the entry
for $B=\{(0,0,0)\}$, $F=\{(5,5,5)\}$ is $(1,1,4)$, and if that cell is occupied as well, the entry for
$F=\{(1,1,4),(5,5,5)\}$ is $(1,4,4)$. The $35$ entries of round $4$ continue in the same way; for instance,
for $B=\{(0,0,0),(1,1,4),(4,4,1)\}$ the moves for $F=\{(5,5,5)\}$, $F=\{(2,3,2),(5,5,5)\}$ and
$F=\{(0,5,5),(2,3,2),(5,0,5),(5,5,5)\}$ are $(2,3,2)$, $(5,0,5)$ and $(0,0,5)$.

\subsubsection*{Two complete plays}
We followed the table from the root in two extreme scenarios, translating every move back to one fixed
frame. (The choice of the symmetry is not unique when the position has a non-trivial stabilizer; other
choices give symmetric plays.)

\emph{Play A: nothing is ever revealed.} Suppose that no threat and no danger pair occurs and that every
cell Breaker wants is empty, so that the play always follows child (i). Then Breaker plays, in rounds $1$
to $8$,
\[
(0,0,0),\ (5,5,5),\ (1,1,4),\ (4,4,1),\ (2,3,2),\ (3,2,3),\ (0,5,5),\ (5,0,0),
\]
i.e.\ four antipodal pairs, one on each space diagonal. After round $8$ Breaker's stones are exactly the set
$\Teight$, and the P-state $\Pst(\Teight,\emptyset,8)$ is the unique leaf with $\theta=8$; in canonical
form it is the entry ``\texttt{8 0,5,61,93,122,154,210,215 -}'' of the leaf file. From round $9$ on
Breaker plays a pairing strategy; $\LEAF(\Teight,\emptyset,8)$ guarantees that a pairing exists whatever
admissible position Maker's eight unrevealed stones occupy (\cref{thm:leaf-implies-C}).

\emph{Play B: every wanted cell is occupied.} Suppose instead that, in every round, the first cell Breaker
wants is occupied by an unrevealed stone of Maker (child (ii)), after which all stones of Maker are
revealed and
the second cell is free. Then Breaker first tries $(0,0,0)$, finds it occupied and plays $(1,1,1)$; in
each of the rounds $2$ to $9$ Breaker first tries the antipode of the stone placed in the round before,
finds it occupied, and plays elsewhere. Breaker's stones are, in order,
\[
(1,1,1),\ (5,0,0),\ (0,0,5),\ (0,5,0),\ (1,4,4),\ (4,4,1),\ (3,3,2),\ (2,3,3),\ (2,2,2),
\]
and the revealed stones of Maker are $(0,0,0)$ and the antipodes of the first eight of them. After round $9$
the state is a leaf with $\theta=0$ (it is the last entry of the leaf file); Breaker's nine stones have
slack $9$, and since all nine stones of Maker are revealed, the leaf condition only requires a pairing for this
one position. So even when Maker ``steals'' every cell that Breaker wants, the strategy reaches a pairing
position one round later than in Play A.

Between these extremes lie the branches with forced events. A threat reveals the five stones of Maker on the
threatening line and a danger pair the eight stones of its two lines, and either costs Breaker a move on a
cell that was not chosen from the table; the certificate then continues from the corresponding child of type (c) or (d).
All these branches end in a leaf after round $10$ at the latest.
